\section{Conclusiones}\label{sec:conclusiones}
Se cumplió el objetivo general de modelar numéricamente el sistema UMZ-SIRS y analizar su estabilidad numérica y la persistencia de la epidemia bajo separación de escalas. En estabilidad numérica, el sistema de seis ecuaciones se formuló e implementó con cuatro integradores y las referencias se validaron hasta $\sim10^{-11}$ (objetivo 1); RK4 mostró orden cuatro y RKF45 igualó su precisión con unas cinco veces menos evaluaciones. Frente a la rigidez (objetivo 2), el paso de los métodos explícitos quedó limitado por el modo rápido ($h<2.04$ y $h<2.82$) con independencia de $\varepsilon$, mientras que Euler hacia atrás conservó $h=4$ con solución directa y quedó limitado a $h\lesssim3.3$ por la convergencia de Gauss--Seidel; a precisión fija, el menor costo lo logró el integrador implícito adaptativo de orden variable. En comportamiento epidemiológico, la variación de parámetros (objetivo 3) mostró que el régimen informativo ($M$ o $Z$) determinó la infección final, con aumentos de hasta 42\,\% y reducciones de hasta 45\,\%, y los mapas (objetivo 4) mostraron que la información indujo endemia con $R_{00}<1$ en el 8\,\% y el 41\,\% de los puntos, mientras el escepticismo extinguió la endemia en el escenario II.

La contribución es un marco numérico reproducible que separa dos tipos de conclusiones: la estabilidad de los esquemas, que depende del método y del solver lineal, y el comportamiento epidemiológico, que depende del régimen informativo y de los umbrales $R_{00}$, $R_{01}$ y $R_{02}$. Su alcance se limita a horizontes finitos, a un conjunto de parámetros exploratorio y a los escenarios de Bulai et al. (2024), sin calibración con datos reales; por ello los resultados describen el comportamiento del modelo y no predicen una epidemia concreta. Queda por estudiar la coexistencia de desinformados y escépticos en los puntos de régimen indefinido.
